% !TEX TS-program = pdflatex
% !TeX program = pdflatex
% !TEX encoding = UTF-8
% !TEX spellcheck = fr

\documentclass[xcolor=table]{beamer}


%\usepackage{fullpage}
%\usepackage[left=2.8cm,right=2.2cm,top=2 cm,bottom=2 cm]{geometry}
\setbeamersize{text margin left=10pt,text margin right=10pt}
\usepackage{amsmath,amssymb} 
\usepackage[T1]{fontenc}
\usepackage{textcomp}
\usepackage[utf8]{inputenc}
\usepackage[french, onelanguage]{algorithm2e}
\usepackage[french]{babel}
\usepackage{arabtex}
\usepackage{txfonts}
\usepackage[]{graphicx}
\usepackage{multirow}
\usepackage{hyperref}
\usepackage{colortbl}
\usepackage{listingsutf8}
\usepackage{wrapfig}
\usepackage{multicol}
\usepackage[export]{adjustbox} %for images in table, also for frame

\hypersetup{
	colorlinks,
	urlcolor = blue
}

%\renewcommand{\baselinestretch}{1.5}

\def\supit#1{\raisebox{0.8ex}{\small\it #1}\hspace{0.05em}}

\AtBeginSection{%
	\begin{frame}
		\sectionpage
	\end{frame}
}

\newcommand{\rottext}[2]{%
	\rotatebox{90}{%
	\begin{minipage}{#1}%
		\raggedleft#2%
	\end{minipage}%
	}%
}

\usepackage{longtable}
\usepackage{tabu}


\institute{ %
École  nationale Supérieure d'Informatique (ESI, ex. INI), Algérie
}
\author[ \textbf{\footnotesize\insertframenumber/\inserttotalframenumber} \hspace*{1.5cm} ESI - BENATCHBA \& ARIES (2019-2020)] %
{BENACHBA Karima \hfill ARIES Abdelkrime}

\titlegraphic{\includegraphics[height=1cm]{assets/esi-logo.png}}%\hspace*{4.75cm}~

\date{Année universitaire : 2019/2020} %\today

\usetheme{Warsaw} % Antibes Boadilla Warsaw

\beamertemplatenavigationsymbolsempty

%\setbeamertemplate{headline}{}

\definecolor{lightblue}{HTML}{D0D2FF}
\definecolor{lightyellow}{HTML}{FFFFAA}
\definecolor{darkblue}{HTML}{0000BB}
\definecolor{olivegreen}{HTML}{006600}
\definecolor{violet}{HTML}{6600CC}

\newcommand{\keyword}[1]{\textcolor{red}{\bfseries\itshape #1}}
\newcommand{\expword}[1]{\textcolor{olivegreen}{#1}}
\newcommand{\optword}[1]{\textcolor{violet}{\bfseries #1}}

\makeatletter
\newcommand\mysphere{%
	\parbox[t]{10pt}{\raisebox{0.2pt}{\beamer@usesphere{item projected}{bigsphere}}}}
\makeatother

%\let\oldtabular\tabular
%\let\endoldtabular\endtabular
%\renewenvironment{tabular}{\rowcolors{2}{white}{lightblue}\oldtabular\rowcolor{blue}}{\endoldtabular}


\NoAutoSpacing %french autospacing after ":"

\def\graphpath{}

\newcommand{\changegraphpath}[1]{\def\graphpath{#1}}


\newcommand{\vgraphpage}[2][.84\textheight]{%
%	\begin{center}%
		\includegraphics[height=#1]{\graphpath #2}%
%	\end{center}%
}

\newcommand{\hgraphpage}[2][\textwidth]{%
%	\begin{center}%
		\includegraphics[width=#1]{\graphpath #2}%
%	\end{center}%
}

\newcommand{\graphpage}[2][]{%
	\includegraphics[#1]{\graphpath #2}%
}

\bibliographystyle{acm}

\newcommand{\insertbibliography}[2]{
	\appendix
	\section*{Bibliographie}
	\nocite{#2}
%	\makeatletter % to change template
%	\setbeamertemplate{headline}[default] % not mandatory, but I though it was better to set it blank
%	\def\beamer@entrycode{\vspace*{-\headheight}} % here is the part we are interested in :)
%	\makeatother
	\begin{multicols*}{2}[\frametitle{\insertsection} \usebeamertemplate*{frametitle}]%\usebeamertemplate*{frametitle}\frametitle{Références}
		\tiny
		\bibliography{#1}
	\end{multicols*}
}

\newcommand{\insertlicence}{
	\begin{frame}[plain]
	\frametitle{Licence : CC-BY-NC 4.0}
%	\framesubtitle{Licence: CC-BY-NC 4.0}
	\begin{center}
		\bfseries Attribution - Pas d'Utilisation Commerciale 4.0 International 
		\url{https://creativecommons.org/licenses/by-nc/4.0/deed.fr}
	\end{center}

	\tiny

	Vous êtes autorisé à : 
	\begin{itemize}
		\item \textbf{Partager} — copier, distribuer et communiquer le matériel par tous moyens et sous tous formats
		\item \textbf{Adapter} — remixer, transformer et créer à partir du matériel
	\end{itemize}
	
	Selon les conditions suivantes : 
	\begin{itemize}
		\item \textbf{Attribution} — Vous devez créditer l'Œuvre, intégrer un lien vers la licence et indiquer si des modifications ont été effectuées à l'Oeuvre. Vous devez indiquer ces informations par tous les moyens raisonnables, sans toutefois suggérer que l'Offrant vous soutient ou soutient la façon dont vous avez utilisé son Oeuvre.
		\item \textbf{Pas d'Utilisation Commerciale} — Vous n'êtes pas autorisé à faire un usage commercial de cette Oeuvre, tout ou partie du matériel la composant. 
		\item \textbf{Pas de restrictions complémentaires} — Vous n'êtes pas autorisé à appliquer des conditions légales ou des mesures techniques qui restreindraient légalement autrui à utiliser l'Oeuvre dans les conditions décrites par la licence.
	\end{itemize}

	\end{frame}
}

\settowidth{\leftmargini}{\usebeamertemplate{itemize item}}
\addtolength{\leftmargini}{\labelsep}

\AtBeginDocument{
	\newcolumntype{L}[2]{>{\vbox to #2\bgroup\vfill\flushleft}p{#1}<{\egroup}} 
	
	\begin{frame}[plain]
		\maketitle
	\end{frame}

%	\insertlicence
}


% needs etoolbox; to break links after -
\appto\UrlBreaks{\do\-}


\title[ML : Arbres de décision (TP)] %
{Machine Learning \\Arbres de décision (support TP)} 

\changegraphpath{../img/arbres/}

\begin{document}
	
\begin{frame}
	\frametitle{Machine Learning}
	\framesubtitle{Arbres de décision (support TP) : Introduction}
	
	\begin{itemize}
		\item Les arbres de décision
		\begin{itemize}
			\item comment créer un arbre de décision en utilisant ID3
			\item en utilisant la régression logistique
			\item dans le cas binaire (appartient à une classe ou hors classe)
		\end{itemize}
	\end{itemize}
\end{frame}


\begin{frame}
	\frametitle{Machine Learning}
	\framesubtitle{Arbres de décision (support TP) : Plan}
	
	\begin{multicols}{2}
		%	\small
		\tableofcontents
	\end{multicols}
\end{frame}

\section{Arbres de décision}

\begin{frame}
	\frametitle{Arbres de décision}
	
	\begin{itemize}
		\item \optword{ID3 (Iterative Dichotomiser 3)} : dévelopé en 1986 par Ross Quinlan. Il peut être appliqué seulement sur les caractéristiques nominales. Il est utilisé pour le classement.
		\item \optword{C4.5} : une extension de ID3 par Ross Quinlan. Il peut être appliqué sur tous les types de caractéristiques. Il est utilisé pour le classement.
		\item \optword{C5.0} : une extension commerciale de C4.5, toujours par Ross Quinlan.
		\item \optword{CART (Classification and Regression Trees)} : comme C4.5 mais utilise d'autres métriques. Aussi, l'algorithme supporte la régression.
	\end{itemize}
	
\end{frame}

\subsection{Algorithmes}

\begin{frame}
	\frametitle{Arbres de décision}
	\framesubtitle{Algorithmes : Construction (apprentissage)}
	
	\begin{algorithm}[H]
		\KwData{$ X, Y$}
		\KwResult{$ (S, A) $ : un arbre} 
		créer un sommet $ s \in S $\;
		
		\eIf{$Y$ est homogène ou on a atteint une condition d'arrêt}{
			$s.classe=\arg\max_{k}{|\{y \in Y / y = k \}|}$ \tcp{$s$ est une feuille}
		}{
			déterminer la caractéristique $X_j$ qui divise le mieux $Y$\;
			diviser $Y$ en sous-ensembles $Y_1, \ldots, Y_K$ selon les valeurs de $X_j$\;
			diviser $X$ en sous-ensembles $X_{*,1}, \ldots, X_{*, K}$ selon les valeurs de $X_j$\;
			$s.j = j$\;
			\ForEach{$ k \in \{1, \ldots, K\}$}{
				$s_k = Construction(X_{*,k}, Y_k) \in S$ $(s, s_k) \in A$ $val(s, s_k) = k$\;
			}
		}
		\caption{Génération d'un arbre de décision (algorithme général)}
	\end{algorithm}
	
\end{frame}

\begin{frame}
	\frametitle{Arbres de décision}
	\framesubtitle{Algorithmes : Recherche (prédiction)}
	
	\begin{algorithm}[H]
		\KwData{$x$ : un échantillon; $(S, A)$ : l'arbre de décision}
		\KwResult{$ y $ : sa classe} 
		
		$s$ est la racine de $(S, A)$\;
		
		
		\eIf{$s$ est une feuille}{
			\Return $s.classe$\;
		}{
			$j = s.j$\; 
			choisir l'arrête $(s, s_k)$ tel que $val(s, s_k) \equiv x_j$\;
			\Return Parcourir(x, $(s_k, A)$)\;
		}
		\caption{Parcours d'un arbre de décision (algorithme général)}
	\end{algorithm}
	
\end{frame}

\subsection{Critique}

\begin{frame}
	\frametitle{Arbres de décision}
	\framesubtitle{Critique : Avantages}
	
	\begin{itemize}
		\item simples à comprendre et à interpréter. On peut visualiser les arbres. Aussi, on peut expliquer les résultats obtenus facilement.
		\item peuvent travailler sur des données avec peu de préparation. Par exemple, ils n'ont pas besoin de la normalisation des données.
		\item acceptent les données numériques et nominales. Les autres algorithmes d'apprentissage sont spécialisés dans un seul type de données.
		\item donnent de bonne performance même si leurs hypothèses sont un peu violées par le modèle réel à partir duquel les données ont été générées.
	\end{itemize}
	
\end{frame}

\begin{frame}
	\frametitle{Arbres de décision}
	\framesubtitle{Critique : Inconvénients}
	
	\begin{itemize}
		\item peuvent être aussi complexes, ils ne généralisent pas bien (overfitting: surapprentissage). On peut régler ça en fixant le nombre minimum des échantillons dans les feuilles ou en fixant la profondeur maximale de l'arbre.
		\item peuvent être instables à cause des variations des données.
		\item il existe des problèmes qui sont un peu difficiles à apprendre par les arbres de décision. Ils ne sont pas faciles à exprimer, par exemple: XOR.
		\item peuvent être biaisés à la classe dominante. Donc, il faut balancer les données avant d'entrainer le système.
		\item ce n'est pas garanti de tomber sur l'arbre de décision optimal.
	\end{itemize}
	
\end{frame}

\section{ID3 (Rappel)}

\begin{frame}
	\frametitle{ID3 (Rappel)}
	
	\begin{itemize}
		\item \keyword{Iterative Dichotomiser 3}
		\item n'accepte que les caractéristiques nominales
		\item seulement pour le classement
		\item on s'arrête si les ensembles des classes dans les feuilles sont homogènes
	\end{itemize}
	
\end{frame}

\subsection{Homogénéité d'un ensemble}

\begin{frame}
	\frametitle{ID3 (Rappel)}
	\framesubtitle{Homogénéité d'un ensemble}
	
	\begin{itemize}
		\item Entropie de Shannon $H(Y)$ pour mesurer l'incertitude d'un ensemble $Y$
		\item $H(Y)=0$ : $Y$ contient les mêmes valeurs (classes)
		\item $H(Y)=1$ : $Y$ contient des valeurs totalement différentes
		\item étant donné $V_y$ l'ensemble du vocabulaire de $Y$ (différentes valeurs ou classes)
	\end{itemize}

	\[H(Y) = - \sum\limits_{v \in V_y} p(v/Y) \log_2 p(v/Y)\]
	
	\[p(v/Y) = \frac{|\{y / y \in Y \wedge y = v\}|}{|S|}\]
	
\end{frame}

\subsection{Division d'un ensemble}

\begin{frame}
	\frametitle{ID3 (Rappel)}
	\framesubtitle{Division d'un ensemble}
	
	\begin{itemize}
		\item $Y$ : un ensemble des prédictions 
		\item $X_j$ : les valeurs de la caractéristique $j$ des échantillons ayant comme prédictions $Y$
		\item $v$ : une valeur parmi les valeurs possibles de $X_j$ (vocabulaire $V_j$)
		\item pour chaque valeur $v \in V_j$, on crée un ensemble $Y_{j, v}$
		\item si $|V| = K$, on aura $K$ ensembles $Y_1, \cdots, Y_K$ et 
	\end{itemize}
	
	\[Y_{j,v} = \{y^{(i)} \in Y / X_j^{(i)} \in X_j \wedge X_j^{(i)} = v\}\]

	
\end{frame}

\subsection{Choix de caractéristique de division}

\begin{frame}
	\frametitle{ID3 (Rappel)}
	\framesubtitle{Choix de caractéristique de division}
	
	\begin{itemize}
		\item On utilise le gain d'entropie (information gain) $IG(Y, X_j)$
		\item $j = \arg\max_{j'} IG(Y, X_{j'})$
		\item Cette mesure calcule combien d'incertitude dans $Y$ a été réduite après sa division en utilisant l'attribut $j$.
		\item c'est la différence entre l'entropie de $Y$ et l'entropie pondérée des ensembles divisées
	\end{itemize}

	\[IG(Y, X_j) = H(Y) - \sum_{v \in V_j} p(v/X_j) H(Y_{j, v})\]
	
	
\end{frame}

\section{CART}

\begin{frame}
	\frametitle{CART}
	
	\begin{itemize}
		\item \keyword{Classification and Regression Trees}
		\item supporte la régression.
		\item essaye de minimiser une fonction de coût.
		\item utilise le pré-élagage en utilisant un critère d'arrêt.
		\item crée des arbres binaires 
	\end{itemize}
	
\end{frame}

\subsection{Homogénéité d'un ensemble}

\begin{frame}
	\frametitle{CART}
	\framesubtitle{Homogénéité d'un ensemble}
	
	\begin{itemize}
		\item indexe de diversité $Gini(Y)$ pour mesurer l'erreur de classification de $Y$
		\item $Gini(Y)=0$ représente la meilleur division
		\item $H(Y)=0.5$ représente la division la plus pire
		\item étant donné $V_y$ l'ensemble du vocabulaire de $Y$ (différentes valeurs ou classes)
		\item Dans le cas de la régression, on utilise \keyword{MSE}
	\end{itemize}
	
	\[Gini(S) = \sum\limits_{v \in V_y} p(v/Y) (1-p(v/Y)) = 1 - \sum\limits_{v \in V_y} p(v/Y)^2 \]
	
	\[p(v/Y) = \frac{|\{y / y \in Y \wedge y = v\}|}{|S|}\]
	
\end{frame}

\subsection{Division d'un ensemble}

\begin{frame}
	\frametitle{CART}
	\framesubtitle{Division d'un ensemble}
	
	\begin{itemize}
		\item $Y$ : un ensemble des prédictions 
		\item $X_j$ : les valeurs de la caractéristique $j$ des échantillons ayant comme prédictions $Y$
		\item $v$ : une valeur parmi les valeurs possibles de $X_j$ (vocabulaire $V_j$)
		\item pour chaque valeur $v \in V_j$, on crée deux ensembles $Y_G$ et $Y_D$
		\item Donc, on aura 2 ensembles $Y_G$ (avec des valeurs $X_j > v$) et $Y_D$ (avec des valeurs $X_j \le v$)
	\end{itemize}
	
	\[Y_G = \{y^{(i)} \in Y / X_j^{(i)} \in X_j \wedge X_j^{(i)} > v\}\]
	\[Y_D = \{y^{(i)} \in Y / X_j^{(i)} \in X_j \wedge X_j^{(i)} \le v\}\]
	
\end{frame}

\subsection{Choix de caractéristique de division}

\begin{frame}
	\frametitle{CART}
	\framesubtitle{Choix de caractéristique de division}
	
	\begin{itemize}
		\item On utilise la diversité Gini de la division $Gini_{div}(Y_G, Y_D)$
		\item $j, v = \arg\min_{j', v'} Gini_{div}(Y_G, Y_D)$
		\item Ici, on ne cherche pas seulement la caractéristique $j$, mais aussi la valeur $v$ qui minimise la diversité Gini de la division
		\item $|S| = |S_G| + |S_D|$
		\item pour la régression, on utilise \keyword{MSE} à la place de \keyword{Gini}
		\item l'estimation dans le cas de la régression est la moyenne des $Y$ dans la feuille
	\end{itemize}
	
	\[Gini_{div}(S_G, S_D) = \frac{|S_G|}{|S|} Gini(S_G) + \frac{|S_D|}{|S|} Gini(S_D)\]
	
	
\end{frame}

%\insertbibliography{ML-preparation}{*}

\end{document}



